\documentclass[11pt,letterpaper]{article}

\usepackage[utf8]{inputenc}
\usepackage[T1]{fontenc}
\usepackage{lmodern}
\usepackage[spanish,es-nodecimaldot]{babel}
\usepackage{geometry}
\usepackage{enumitem}
\usepackage{xcolor}
\usepackage{hyperref}
\usepackage{titlesec}
\usepackage{tabularx}
\usepackage{microtype}
\usepackage{needspace}

\geometry{
  top=1.5cm,
  bottom=1.5cm,
  left=1.8cm,
  right=1.8cm
}

% Paleta tomada del portfolio: verde bosque, cuero cognac y crema cálido.
\definecolor{forest}{HTML}{173A32}
\definecolor{sage}{HTML}{497466}
\definecolor{cognac}{HTML}{986846}
\definecolor{gold}{HTML}{BDAA79}
\definecolor{cream}{HTML}{F5F0E7}
\definecolor{accent}{HTML}{8A5A3A}
\definecolor{accentdark}{HTML}{173A32}
\definecolor{accentlight}{HTML}{D9E3DD}
\definecolor{bodytext}{HTML}{292D2B}
\definecolor{muted}{HTML}{68736E}

\hypersetup{
  colorlinks=true,
  urlcolor=accent,
  linkcolor=accent,
  pdftitle={CV de Nahuel Mosse},
  pdfauthor={Nahuel Mosse},
  pdfsubject={Technical Lead y Software Engineer},
  pdfkeywords={Technical Lead, Senior Full Stack Developer, TypeScript, Next.js, NestJS, Azure,
    Terraform, Kubernetes, OpenAI, arquitectura de software, liderazgo técnico}
}

\pagestyle{empty}
\renewcommand{\familydefault}{\sfdefault}
\setlength{\parindent}{0pt}
\setlength{\parskip}{0.25em}
\setlist[itemize]{
  leftmargin=1.25em,
  itemsep=0.12em,
  topsep=0.2em,
  parsep=0pt,
  partopsep=0pt
}
\setlist[itemize,2]{leftmargin=1.4em,label=--}
\urlstyle{same}
\color{bodytext}

\newcommand{\cvsection}[1]{%
  \Needspace{4\baselineskip}%
  \vspace{0.55em}%
  \noindent
  {\color{accent}\rule[-0.18em]{3pt}{1.25em}}%
  \hspace{0.45em}%
  {\large\bfseries\color{accentdark}\MakeUppercase{#1}}%
  \hspace{0.55em}%
  {\color{accentlight}\leaders\hrule height 1.2pt\hfill\kern0pt}\par
  \vspace{0.28em}%
}

\newcommand{\cventry}[3]{%
  \begin{tabularx}{\textwidth}{@{}Xr@{}}
    \textbf{\color{accentdark}#1} \textcolor{muted}{\textbar}\enspace #2
    & \small\textcolor{muted}{#3}
  \end{tabularx}
}

\newcommand{\project}[1]{%
  \Needspace{6\baselineskip}%
  \vspace{0.3em}%
  {\color{accent}\rule[0.18em]{0.45em}{0.45em}}\hspace{0.4em}%
  \textbf{\color{accentdark}Proyecto: #1}\par
}

\begin{document}

\noindent
\colorbox{forest}{%
  \parbox{\dimexpr\textwidth-2\fboxsep\relax}{%
    \centering
    \vspace{0.7em}
    {\fontsize{27}{30}\selectfont\bfseries\color{cream}Nahuel Mosse}\par
    \vspace{0.12em}
    {\Large\color{cream}Technical Lead}\par
    \vspace{0.18em}
    {\scriptsize\ttfamily\color{gold}%
      \textless from curiosity to craft /\textgreater
    }\par
    \vspace{0.48em}
    {\small\color{cream}%
      \begingroup\hypersetup{urlcolor=cream}%
      \href{mailto:contact@nahuelmosse.com}{contact@nahuelmosse.com}
      \enspace\textbar\enspace
      Buenos Aires, Argentina (AMBA)\par
      \vspace{0.12em}
      \href{https://linkedin.com/in/nahuelmosse}{linkedin.com/in/nahuelmosse}
      \enspace\textbar\enspace
      \href{https://github.com/NahuelMosse}{github.com/NahuelMosse}
      \enspace\textbar\enspace
      \href{https://gitlab.com/NahuelMosse}{gitlab.com/NahuelMosse}%
      \endgroup
    }\par
    \vspace{0.6em}
  }%
}
{\color{cognac}\rule{\textwidth}{1.4pt}}

\cvsection{Perfil profesional}

Technical Lead con más de 4 años liderando productos web end-to-end. Integra
arquitectura, desarrollo full stack, business analysis, liderazgo técnico y DevOps/cloud
para convertir necesidades de negocio en plataformas escalables, seguras y mantenibles.
Especializado en APIs e IA aplicada con impacto medible.

\cvsection{Experiencia laboral}

\cventry{Consultora Kovix}{Senior Full Stack Developer / Technical Lead}{06/2022 -- Actualidad}

\project{Smart Warehousing --- AI Engineering Enablement Lead / SDD Coach
  \textnormal{\small\textcolor{muted}{\enspace 09/2026 -- Actualidad}}}
\begin{itemize}
  \item Creación y evolución de una metodología transversal de desarrollo asistido por IA,
    aplicable a proyectos greenfield y brownfield de diferentes equipos y arquitecturas.
  \item Capacitación de empleados mediante formación, documentación y acompañamiento
    práctico para adoptar SDD y Claude Code multiagente en su trabajo cotidiano.
  \item Facilitación de la adopción con un rol análogo al de Scrum Master, enfocado en SDD:
    asesoramiento a los equipos, resolución de bloqueos, seguimiento del uso de la
    metodología y mejora continua a partir del feedback.
  \item Definición de workflows con especificaciones, responsabilidades, criterios de
    aceptación y validaciones explícitas para delegar implementación a agentes de IA bajo
    revisión y gobierno humano.
  \item Modelo operativo donde los agentes asumen tareas de alcance junior/semi-senior y
    los perfiles Senior concentran arquitectura, escalabilidad, mantenibilidad, calidad y
    decisiones críticas del producto.
\end{itemize}

\project{Grupo ST --- Plataforma de Riesgo Crediticio
  \textnormal{\small\textcolor{muted}{\enspace 08/2026 -- Actualidad}}}
\begin{itemize}
  \item Technical Lead en asignación paralela con Visma, dirigiendo una plataforma de
    gestión documental y decisiones orientada a automatizar la aprobación de créditos.
  \item Modelo de desarrollo AI-first: orquestación de agentes con Claude Code y
    Spec-Driven Development (SDD) para implementar el producto end-to-end desde
    especificaciones y decisiones de arquitectura.
  \item Diseño de una arquitectura DDD sobre Next.js y NestJS, con Sequelize y PostgreSQL;
    definición de estándares técnicos y control de la evolución del producto aún en desarrollo.
  \item Estrategia de calidad para desarrollo con IA mediante Jest y una suite de regresión
    end-to-end con Playwright, validando los flujos completos para prevenir regresiones.
  \item \textbf{Stack:} TypeScript, Next.js, React, NestJS, Sequelize, PostgreSQL, DDD,
    Claude Code, SDD, Jest, Playwright y Turborepo.
\end{itemize}

\newpage
\project{Visma Payroll Onboarding \textnormal{\small\textcolor{muted}{\enspace 01/2026 -- Actualidad}}}
\begin{itemize}
  \item Liderazgo técnico y elicitación de requisitos: interlocutor directo con el cliente,
    convirtiendo necesidades del negocio de payroll en alcance, arquitectura y roadmap.
  \item Gestión del calendario, prioridades y distribución del trabajo; mentoría y
    capacitación de un desarrollador junior, y conducción del design system del producto.
  \item Diseño e implementación en un monorepo Nx con Next.js y NestJS, aplicando Clean
    Architecture y DDD; integración con APIs internas y Azure OpenAI que redujo una etapa
    del onboarding de 2 meses a 1 semana.
  \item Arquitectura de datos y procesos asíncronos con Prisma/PostgreSQL para persistencia
    y BullMQ/Redis para colas, caché y rate limiting distribuido de APIs e IA; optimización
    de costos evitando usar Redis persistente como almacenamiento durable.
  \item Responsable de ciberseguridad, CI/CD e infraestructura como código: despliegue en
    Azure con Terraform y AKS, identidades administradas, Key Vault, redes privadas y
    políticas de seguridad sobre Kubernetes.
  \item \textbf{Aplicación y datos:} TypeScript, Nx, Next.js, NestJS,
    Prisma/PostgreSQL, Redis/BullMQ y Azure OpenAI.
  \item \textbf{Cloud y calidad:} Terraform, AKS/Kubernetes, Azure DevOps, Playwright y
    Claude Code con flujos multiagente.
\end{itemize}

\project{Frávega --- WMS B2C \textnormal{\small\textcolor{muted}{\enspace 06/2023 -- 01/2026}}}
\begin{itemize}
  \item Modernización del fulfillment de pedidos en un warehouse de gran escala:
    desarrollo del software que orquestaba el armado de paquetes y su integración con el
    sorter responsable del transporte interno.
  \item Implementación integral de la arquitectura BFF del módulo Sellers, soportando
    40.000 requests diarias (1,2 millones mensuales).
  \item Simplificación de procesos de negocio en un 80\% y reducción del tiempo promedio
    de recepción de 7 h 48 min a 46 min (aprox. 90\%).
  \item Diseño técnico y documentación de contratos entre APIs, optimización de queries,
    implementación de unit testing y soporte en producción.
  \item Planificación estratégica de releases críticos, gestión preventiva de riesgos de
    calendario, elaboración de guías de despliegue y mentoría del equipo.
  \item \textbf{Stack:} Node.js, NestJS, Angular, React, Jest, Vitest, TypeScript,
    Sequelize, SQL, PostgreSQL, Tailwind CSS, HTML, CSS y webhooks.
\end{itemize}

\project{Jüsto --- WMS Multiwarehouse
  \textnormal{\small\textcolor{muted}{\enspace 11/2022 -- 12/2023}}}
\begin{itemize}
  \item Diseño de un WMS multiwarehouse con aislamiento similar a multitenancy para los
    centros de distribución, más una versión reducida para Micro-Fulfillment Centers (MFC).
  \item Sustitución de planillas y conocimiento informal por trazabilidad de inventario y
    ubicaciones, con despacho FEFO (First Expired, First Out) para el supermercado online.
  \item Más de 10 despliegues, 50\% menos tiempo de preparación, 40 millones de operaciones
    y mayor holgura ante picos de ventas.
  \item Diseño técnico y contratos entre APIs; estimación de esfuerzos, asignación y
    seguimiento de tareas, coordinación del equipo e implementación de unit testing.
  \item Mentoría y colaboración con el equipo para resolver bloqueos y sostener la
    estabilidad durante los picos históricos de ventas.
  \item \textbf{Stack:} Node.js, NestJS, Angular, React, Jest, TypeScript,
    Sequelize, SQL, PostgreSQL, Tailwind CSS, HTML, CSS y webhooks.
\end{itemize}

\newpage
\cvsection{Educación}

\cventry{Ingeniería Informática}{Universidad de Morón}{03/2022 -- Actualidad}
\begin{itemize}
  \item Cursando 4.º año.
  \item Beca al mérito con promedio superior a 8,5.
  \item Líder de grupos de estudio.
\end{itemize}

\cventry{Analista en Informática}{Universidad de Morón}{03/2022 -- 12/2025}
\begin{itemize}
  \item Carrera finalizada.
  \item Beca al mérito con promedio superior a 8,5.
  \item Líder de grupos de estudio.
\end{itemize}

\cventry{Técnico Electromecánico}{Colegio Emaús}{03/2015 -- 12/2021}
\begin{itemize}
  \item Orientación en electromecánica.
\end{itemize}

\cvsection{Proyectos personales}

\cventry{Nasa Calendar}{Next.js}{}

Calendario que consume una API de la NASA para mostrar la imagen astronómica de cada día.

\href{https://github.com/NahuelMosse/nasa-calendar}{github.com/NahuelMosse/nasa-calendar}

\begin{itemize}
  \item \textbf{Stack:} TypeScript, JavaScript, server-side rendering (SSR), React.js,
    Next.js y Tailwind CSS.
\end{itemize}

\cvsection{Logros académicos y certificaciones}

{\large\bfseries\color{accentdark}Certificaciones y formación complementaria}\par
\begin{itemize}
  \item \textbf{EF SET English Certificate --- C2 Proficient} --- EF Education First
    (2026). Certificación del nivel más alto del MCER.
  \item \textbf{Curso de Arquitectura de Software Aplicada} --- Platzi (2026).
    Patrones avanzados de arquitectura y diseño de sistemas.
  \item \textbf{Curso de Fundamentos de Arquitectura de Software} --- Platzi (2026).
    Principios y buenas prácticas de arquitectura.
  \item \textbf{Liceo Cultural Británico --- Nivel 8}. Inglés avanzado y
    literatura inglesa.
  \item \textbf{Liceo Cultural Británico --- Nivel 6} (2025). Formación orientada a la
    preparación del First Certificate.
  \item \textbf{Fundamentos de Arquitectura de Software} --- Platzi (2023).
  \item \textbf{Fundamentos de Ingeniería de Software} --- Platzi (2023).
  \item \textbf{Responsive Web Design} --- freeCodeCamp (2022). HTML5, CSS3, Flexbox y
    CSS Grid.
\end{itemize}

\newpage
{\large\bfseries\color{accentdark}Olimpíadas y competencias académicas}\par
\begin{itemize}
  \item \textbf{Finalista nacional --- Olimpiada Argentina de Física}, modalidad Juegos
    (OAF, 2019).
  \item \textbf{Participante de la instancia metropolitana --- Olimpiada Matemática
    Argentina} (OMA, 2019).
  \item \textbf{Campeón de Desafíos Científicos --- Ematemática}, Colegio Emaús (2019).
  \item \textbf{Campeón de la tercera instancia --- Ematemática}, Colegio Emaús (2019).
  \item \textbf{Campeón de la segunda instancia --- Ematemática}, Colegio Emaús (2019).
  \item \textbf{Subcampeón de la primera instancia --- Ematemática}, Colegio Emaús (2019).
  \item \textbf{Segundo subcampeón nacional --- Mateclubes}, Olimpiada Matemática
    Argentina (OMA, 2018).
  \item \textbf{Finalista nacional --- Olimpiada de Literatura y Matemática} (2018).
  \item \textbf{Participante de la instancia metropolitana --- Olimpiada Matemática
    Argentina} (OMA, 2018).
  \item \textbf{Campeón de Desafíos Científicos --- Ematemática}, Colegio Emaús (2018).
  \item \textbf{Segundo subcampeón de la tercera instancia --- Ematemática}, Colegio
    Emaús (2018).
  \item \textbf{Campeón de la segunda instancia --- Ematemática}, Colegio Emaús (2018).
  \item \textbf{Segundo subcampeón de la primera instancia --- Ematemática}, Colegio
    Emaús (2018).
  \item \textbf{Finalista nacional --- Mateclubes}, Olimpiada Matemática Argentina
    (OMA, 2017).
  \item \textbf{Campeón --- Ematemática}, Colegio Emaús (2017).
\end{itemize}

\newpage
\cvsection{Habilidades técnicas}

\begin{itemize}
  \item \textbf{Lenguajes:} TypeScript, JavaScript, SQL, Python, Go, C++ y Java.
  \item \textbf{Frontend/UI:} React, Next.js, Angular, Tailwind CSS, shadcn/ui, Radix UI,
    SSR y design systems.
  \item \textbf{Backend y datos:} Node.js, NestJS, Prisma, Sequelize, PostgreSQL, MySQL,
    MongoDB, Redis, BullMQ, REST APIs y webhooks.
  \item \textbf{Arquitectura:} DDD, Clean Architecture, BFF, sistemas distribuidos,
    contratos entre APIs y procesamiento asíncrono.
  \item \textbf{Cloud, DevOps y seguridad:} Azure, AKS, Kubernetes, Terraform, Docker,
    Azure DevOps, CI/CD, Key Vault, identidades administradas y redes privadas.
  \item \textbf{AI Engineering:} OpenAI, Azure OpenAI, Claude Code, orquestación
    multiagente, SDD y workflows de desarrollo asistido por IA.
  \item \textbf{Testing y observabilidad:} Jest, Vitest, Playwright, Datadog, OpenSearch
    y Argo.
  \item \textbf{Gestión:} liderazgo técnico, business analysis, elicitación de requisitos,
    planificación, Scrum, Agile, Kanban, Jira y Confluence.
\end{itemize}

\cvsection{Idiomas}

\begin{itemize}
  \item \textbf{Español:} nativo.
  \item \textbf{Inglés:} C2 Proficient (EF SET, 2026).
\end{itemize}

\cvsection{Habilidades blandas}

Liderazgo, planificación, comunicación clara, negociación, resolución de problemas,
aprendizaje continuo y autónomo, trabajo en equipo y manejo de personal.

\cvsection{Recomendaciones}

\textbf{Resumen de LinkedIn}
\begin{itemize}
  \item \textbf{Braian Lo Surdo --- Técnico en Programación Informática:} destaca rapidez,
    eficiencia, capacidad colaborativa, entusiasmo y calidez en el trabajo en equipo.
  \item \textbf{Enzo Codini --- Full Stack Developer:} resalta el aporte técnico y humano
    en proyectos complejos, la buena disposición y la calidad de trabajo, considerándolo
    un referente en el equipo.
\end{itemize}

\end{document}
